% Tài liệu kỹ thuật (portfolio/CV) đi kèm docs/TECHNICAL_OVERVIEW.md.
% Cần XeLaTeX/LuaLaTeX (fontspec + polyglossia) để hiển thị tiếng Việt.
% Biên dịch: tectonic TECHNICAL_OVERVIEW.tex (hoặc xelatex, chạy 2 lần cho
% mục lục / danh sách hình).
\documentclass[11pt,a4paper]{article}

\usepackage{fontspec}
\usepackage{polyglossia}
\setmainlanguage{vietnamese}
\setmainfont{Times New Roman} % đổi sang font có sẵn trên máy nếu không có

\usepackage{geometry}
\geometry{margin=2.3cm}
\usepackage{microtype}
\sloppy
\emergencystretch=3em
\usepackage{booktabs}
\usepackage{hyperref}
\usepackage{xcolor}
\usepackage{enumitem}
\usepackage{listings}

\usepackage{tikz}
\usetikzlibrary{arrows.meta, positioning, shapes.geometric, calc, fit, backgrounds}
\usepackage{pgfplots}
\pgfplotsset{compat=1.18}

\lstset{
  basicstyle=\ttfamily\small,
  breaklines=true,
  frame=single,
  columns=fullflexible,
  backgroundcolor=\color{black!3},
}

\hypersetup{
  colorlinks=true,
  linkcolor=blue!50!black,
  citecolor=blue!50!black,
  urlcolor=blue!50!black,
}

% ---- các style node TikZ dùng lại nhiều lần -------------------------------
\tikzset{
  procbox/.style={rectangle, rounded corners=3pt, draw=blue!55!black, very thick,
    fill=blue!6, text width=3.3cm, minimum height=1.1cm, align=center,
    font=\footnotesize, inner sep=5pt},
  llmbox/.style={rectangle, rounded corners=3pt, draw=orange!65!black, very thick,
    fill=orange!8, text width=3.3cm, minimum height=1.0cm, align=center,
    font=\footnotesize, inner sep=5pt},
  iobox/.style={rectangle, draw=black!70, very thick, fill=black!4,
    text width=3.6cm, minimum height=1.0cm, align=center, font=\footnotesize,
    inner sep=5pt},
  goodbox/.style={rectangle, rounded corners=3pt, draw=green!45!black, very thick,
    fill=green!8, text width=3.4cm, minimum height=1.1cm, align=center,
    font=\footnotesize, inner sep=5pt},
  badbox/.style={rectangle, rounded corners=3pt, draw=red!55!black, very thick,
    fill=red!7, text width=3.4cm, minimum height=1.1cm, align=center,
    font=\footnotesize, inner sep=5pt},
  arr/.style={-{Latex[length=2.3mm]}, thick},
}

\title{\textbf{AI-Powered Automated Grading}\\[2pt]
  \large Tổng quan kỹ thuật \& các quyết định thiết kế}
\author{Khang Nguyen}
\date{\today}

\begin{document}

\maketitle

\begin{abstract}
\noindent Tài liệu này trình bày các quyết định kỹ thuật đứng sau một hệ
thống chấm điểm tự động bài tập lập trình dùng LLM dưới 7 tỷ tham số, xây
dựng quanh một yêu cầu cứng, có thể kiểm chứng: \emph{điểm số khớp tuyệt
đối (exact-match) giữa các lần chạy lặp lại trên cùng một đầu vào}, kể cả
dưới điều kiện serving có batching thật. Đây là bản mở rộng có sơ đồ của
\texttt{docs/TECHNICAL\_OVERVIEW.md} trong repo dự án, viết cho mục đích
portfolio/CV. Phần khung học thuật của cùng dự án (tổng quan tài liệu, mô
tả phương pháp chính thức, kết quả kiểu luận văn) nằm riêng ở
\texttt{docs/thesis/}.
\end{abstract}

\tableofcontents
\listoffigures

% =====================================================================
\section{Vấn đề kỹ thuật thực sự}

``Chấm code bằng LLM'' tự bản thân nó không phải bài toán kỹ thuật khó —
một prompt đơn giản làm được việc đó trong một buổi chiều. Bài toán mà dự
án này thực sự giải quyết là:

\begin{quote}
\itshape Với cùng một đề bài, cùng một rubric, và cùng một bài nộp, hệ
thống có thể được thiết kế để luôn trả về \textbf{đúng cùng một điểm số ở
mọi lần chạy}, kể cả dưới một cấu hình serving có batching như trong sản
xuất thật — và tuyên bố đó có thể được chứng minh bằng một phép đo, chứ
không phải một cờ cấu hình?
\end{quote}

Đó là bài toán tính xác định (determinism) dưới điều kiện serving thật.
Nó cắt ngang qua ba lớp thường không được xem là cùng một bài toán: cấu
trúc đầu ra của mô hình, chiến lược giải mã (decoding), và cơ chế batching
nội bộ của serving engine.

% =====================================================================
\section{Vì sao ``cứ hỏi LLM chấm điểm'' không đáp ứng được yêu cầu này}

Một lệnh gọi tự do (free-form) yêu cầu mô hình trả về một điểm toàn thể
0--100 duy nhất có phương sai giữa các lần chạy rất cao: kể cả ở
\texttt{temperature=0}, mô hình vẫn có nhiều đường suy luận văn bản tự do
khác nhau nhưng đều hợp lý để đi đến cùng một kết luận, và đường nào được
chọn lại nhạy cảm với nhiễu ở tầng serving. Dự án đo trực tiếp khoảng cách
này thay vì tin theo cảm tính. Thực nghiệm ablation
(\texttt{experiments/ablation.py}) so sánh ba kiến trúc trên cùng bộ dữ
liệu 65 mẫu, chấm theo độ tương quan hạng Spearman với nhãn tham chiếu của
bộ dữ liệu.

\begin{figure}[htbp]
\centering
\begin{tikzpicture}
\begin{axis}[
    ybar,
    bar width=26pt,
    width=11cm, height=6.2cm,
    symbolic x coords={freeform, hybridfree, hybriddec},
    xtick=data,
    xticklabels={free\_form, hybrid\_freeform, hybrid\_decomposed},
    xticklabel style={font=\small, align=center},
    ymin=0, ymax=1,
    ylabel={Spearman $\rho$ so với nhãn tham chiếu},
    nodes near coords,
    nodes near coords style={font=\small},
    enlarge x limits=0.3,
    axis lines*=left,
    ymajorgrids,
    major grid style={draw=black!15},
  ]
\addplot[fill=blue!35, draw=blue!55!black] coordinates
  {(freeform,0.60) (hybridfree,0.85) (hybriddec,0.88)};
\end{axis}
\end{tikzpicture}
\caption{Mức độ đồng thuận với nhãn tham chiếu tăng đơn điệu theo từng bước
phân rã: một lệnh gọi free-form
$\rightarrow$ correctness từ sandbox + một lệnh gọi free-form cho phần còn
lại
$\rightarrow$ correctness từ sandbox + một lệnh gọi có cấu trúc cho mỗi
tiêu chí rubric. Bảng đầy đủ kèm p-value trong \texttt{docs/thesis/results.md}.}
\label{fig:ablation-bar}
\end{figure}

Hình~\ref{fig:ablation-bar} là căn cứ thực nghiệm cho kiến trúc ở mục sau —
một kết quả đo được, không phải một sở thích thiết kế.

% =====================================================================
\section{Kiến trúc}

\texttt{src/grader/pipeline.py} là một hàm tuần tự duy nhất — cố tình
\textbf{không} xây trên một agent framework (LangChain, CrewAI, v.v.).
Không có khám phá nhiều bước hay định tuyến công cụ nào cần điều phối ở
đây: chỉ có ba lệnh gọi phân loại độc lập và một phép cộng số học. Mỗi lớp
điều phối thêm vào là một nơi có thể làm rò rỉ thêm phi-xác-định hoặc
token, mà không mang lại lợi ích gì cho khối lượng công việc này.

\begin{figure}[htbp]
\centering
\begin{tikzpicture}[>=Latex]

\node[iobox] (input) at (5.5,14.2) {\textbf{Đầu vào}\\statement, rubric,\\submission\_code, test\_cases};

\node[procbox] (sandbox) at (0,10.4) {\texttt{sandbox.}\\\texttt{run\_submission()}\\correctness\\\emph{(xác định, không LLM)}};

\node[llmbox] (eff)   at (9.4,12.0) {\texttt{check\_efficiency()}\\JSON có cấu trúc, temp=0};
\node[llmbox] (style) at (9.4,9.4) {\texttt{check\_code\_style()}\\JSON có cấu trúc, temp=0};
\node[llmbox] (edge)  at (9.4,6.6) {\texttt{check\_edge\_case\_}\\\texttt{handling()}\\JSON có cấu trúc, temp=0};

\node[procbox] (agg) at (5.5,4.4) {\texttt{aggregator.}\\\texttt{aggregate()}\\tổng có trọng số cố định\\\emph{(code thuần, không LLM)}};

\node[iobox] (output) at (5.5,1.7) {\textbf{final\_score\_0\_to\_100}};

\draw[arr] (input.south west) -- (sandbox.north);
\draw[arr] (input.south) -- (eff.north west);
\draw[arr] (input.south east) -- (style.north west);
\draw[arr] (input.east) -- (edge.north west);

\draw[arr] (sandbox.south) -- (agg.north west);
\draw[arr] (eff.south west) -- (agg.north east);
\draw[arr] (style.west) -- (agg.east);
\draw[arr] (edge.north west) -- (agg.south east);

\draw[arr] (agg.south) -- (output.north);

\end{tikzpicture}
\caption{Toàn bộ pipeline chấm điểm. Correctness chạy trên đường thực thi
thật, không có LLM tham gia; mọi tiêu chí rubric còn lại là một lệnh gọi
LLM hẹp, ràng buộc theo schema; điểm cuối cùng được tính bằng Python thuần,
không bao giờ do mô hình tính.}
\label{fig:pipeline}
\end{figure}

% =====================================================================
\section{Các quyết định kỹ thuật cốt lõi}

\subsection{Correctness không bao giờ chạm vào LLM}

\texttt{src/grader/sandbox.py} chạy code nộp bài với các test case thật
trong một namespace riêng biệt và tính \texttt{pass\_rate} trực tiếp —
không có suy luận mô hình nào trên con đường quan trọng nhất đối với tính
xác định. Điều này cũng làm lộ ra một bài học tinh tế hơn trong quá trình
xây bộ dữ liệu (\texttt{data/script.md}): hai bài nộp \emph{được thiết kế}
để chứa lỗi vẫn đạt 100\% test, vì bộ test không tình cờ bao phủ đúng dạng
lỗi đó. Correctness từ sandbox chỉ đáng tin cậy bằng bộ test đứng sau nó.

\subsection{Giải mã ràng buộc theo schema, không phải ``xin JSON một cách lịch sự''}

Mỗi tiêu chí rubric ngoài correctness có một prompt hẹp riêng
(\texttt{prompts/*.txt}) và một schema Pydantic nhỏ riêng
(\texttt{src/grader/schemas.py}): một \texttt{EfficiencyCheck} (boolean +
một cụm từ ngắn mô tả mẫu quan sát được), một \texttt{StyleCheck} (một
enum thứ bậc 3 giá trị), một \texttt{EdgeCaseCheck} (boolean + danh sách
trường hợp thiếu, có thể rỗng). Các ràng buộc này được áp dụng ngay ở tầng
giải mã qua \texttt{response\_format: json\_schema}, được biên dịch thành
văn phạm GBNF bởi \texttt{llama.cpp}
(\texttt{src/grader/llm\_client.py}) — mô hình \emph{không có khả năng cấu
trúc} để sinh ra bất cứ thứ gì ngoài schema, chứ không chỉ đơn thuần được
prompt để làm vậy.

\begin{figure}[htbp]
\centering
\begin{tikzpicture}[>=Latex]

% -- đường không ràng buộc (trên) --
\node[iobox] (p1) at (0,6) {Prompt: ``chấm bài này,\\trả về JSON''};
\node[llmbox] (m1) at (4.2,6) {Giải mã LLM\\\emph{không ràng buộc}};
\node[badbox] (t1) at (8.6,6) {Văn bản tự do\\(có thể không\\phải JSON hợp lệ)};
\node[badbox] (r1) at (12.6,6) {Bước regex/parse\\\textbf{rủi ro: parse fail,}\\\textbf{trôi giữa các lần chạy}};

\draw[arr] (p1) -- (m1);
\draw[arr] (m1) -- (t1);
\draw[arr] (t1) -- (r1);

% -- đường ràng buộc (dưới) --
\node[iobox] (p2) at (0,2.6) {Prompt: chỉ một\\tiêu chí rubric};
\node[llmbox] (m2) at (4.2,2.6) {Giải mã LLM\\\emph{ràng buộc bởi văn phạm}\\\emph{GBNF (từ schema)}};
\node[goodbox] (t2) at (8.6,2.6) {JSON hợp lệ về\\cấu trúc, khớp\\schema Pydantic};
\node[goodbox] (r2) at (12.6,2.6) {\texttt{schema.model\_validate()}\\\textbf{luôn thành công,}\\\textbf{không có bề mặt trôi}};

\draw[arr] (p2) -- (m2);
\draw[arr] (m2) -- (t2);
\draw[arr] (t2) -- (r2);

\end{tikzpicture}
\caption{Giải mã tự do (trên) để lại một bề mặt mở cho lỗi parse và
phương sai điểm số. Giải mã ràng buộc theo văn phạm với một schema nhỏ
(dưới) loại bỏ bề mặt đó về mặt cấu trúc — đây là cơ chế đứng sau phương
sai bằng 0 trên từng check trong thực nghiệm consistency
(Hình~\ref{fig:consistency}), không chỉ là ý định thiết kế.}
\label{fig:constrained-decoding}
\end{figure}

\subsection{Một lỗi trong công thức ánh xạ điểm đáng biết}

Các điểm thành phần là hàm xác định của đầu ra có cấu trúc, ví dụ:

\begin{lstlisting}[language=Python]
# edge_case_handling: handled=False must never silently score 1.0
score = max(0.25, 1.0 - 0.25 * max(1, len(result.missing_cases)))
\end{lstlisting}

\texttt{max(1, ...)} tồn tại vì một phiên bản trước của công thức này tính
\texttt{1.0 - 0.25 * 0 == 1.0} bất cứ khi nào mô hình trả về
\texttt{handled=False} nhưng danh sách \texttt{missing\_cases} rỗng —
âm thầm ghi đè phán quyết ``không xử lý được'' của chính mô hình bằng một
điểm tuyệt đối. Danh sách rỗng ở đây có nghĩa ``chưa xác định'', không
phải ``không bị trừ điểm''. Đây là một lỗi nhỏ, nhưng đúng kiểu lỗi mà một
lượt kiểm tra nhanh trên happy-path sẽ bỏ sót, và nó trực tiếp làm suy yếu
tuyên bố correctness trung tâm của dự án nếu để sót lại.

\subsection{Tính xác định được đo, không giả định}

\texttt{temperature=0} + seed cố định + giải mã ràng buộc \textbf{không}
đảm bảo lặp lại chính xác từng bit dưới một serving engine có batching —
tính không kết hợp của số học dấu phẩy động qua các tổ hợp batch khác nhau
là một hiện tượng có thật, đã được ghi nhận trong serving LLM, không phải
mối lo lý thuyết. \texttt{experiments/consistency\_test.py} đo trực tiếp
điều này thay vì khẳng định cấu hình làm cho hệ thống có tính xác định.

\begin{figure}[htbp]
\centering
\begin{tikzpicture}[>=Latex]

\node[iobox] (sample) at (5.5,7.6) {Mẫu đầu vào cố định\\(statement + rubric\\+ submission\_code)};

\node[procbox] (seq) at (1.6,5.2) {\texttt{--parallel 1}\\yêu cầu tuần tự\\(không batching)};
\node[procbox] (batch) at (9.4,5.2) {\texttt{--parallel 4}\\8 yêu cầu đồng thời\\(batching thật, xác nhận\\qua slot ID của server)};

\node[llmbox] (run1) at (1.6,3.0) {Toàn bộ pipeline\\$\times$ 100 lần};
\node[llmbox] (run2) at (9.4,3.0) {Toàn bộ pipeline\\$\times$ 100 lần};

\node[goodbox] (res1) at (1.6,0.7) {Số giá trị điểm\\khác nhau: \textbf{1}\\exact-match = 1.0};
\node[goodbox] (res2) at (9.4,0.7) {Số giá trị điểm\\khác nhau: \textbf{1}\\exact-match = 1.0};

\draw[arr] (sample) -- (seq);
\draw[arr] (sample) -- (batch);
\draw[arr] (seq) -- (run1);
\draw[arr] (batch) -- (run2);
\draw[arr] (run1) -- (res1);
\draw[arr] (run2) -- (res2);

\end{tikzpicture}
\caption{Phương pháp đo consistency: cùng một mẫu cố định được chấm 100
lần dưới cả serving không batching và có batching. Kết quả trên 4 mẫu dễ
và 4 mẫu cố tình khó (có lỗi / trường hợp biên): tỷ lệ exact-match = 1.0 ở
mọi điều kiện, phương sai bằng 0 trên mọi breakdown theo từng check. Xem
\texttt{docs/thesis/results.md} để biết các lưu ý về phạm vi (cơ chế
batching cụ thể của llama.cpp, một mức concurrency).}
\label{fig:consistency}
\end{figure}

\subsection{Triển khai xoay quanh giới hạn phần cứng, không bất chấp nó}

Kế hoạch ban đầu giả định dùng vLLM, serving engine phổ biến nhất trong
tài liệu LLM-as-judge mà dự án này dựa vào. Máy phát triển là Apple Silicon
(M1 Pro) không có GPU hỗ trợ CUDA — vLLM không có backend GPU nào được hỗ
trợ trên nền tảng này.

\begin{figure}[htbp]
\centering
\begin{tikzpicture}[>=Latex]

\node[iobox] (start) at (5.2,5.6) {Máy phát triển:\\Apple Silicon (M1 Pro),\\không có GPU CUDA};

\node[badbox] (vllm) at (1.4,3.4) {Kế hoạch ban đầu:\\vLLM\\\textbf{không có backend}\\\textbf{GPU nào phù hợp}};
\node[goodbox] (llama) at (9.0,3.4) {Đã chọn: llama.cpp\\\texttt{llama-server}\\Tăng tốc bằng Metal,\\tương thích API OpenAI};

\node[goodbox] (gbnf) at (9.0,1.0) {Cho phép giải mã\\ràng buộc văn phạm\\GBNF\\$\rightarrow$ Hình~\ref{fig:constrained-decoding}};

\draw[arr] (start) -- (vllm) node[midway, above, sloped, font=\scriptsize] {bị chặn};
\draw[arr] (start) -- (llama) node[midway, above, sloped, font=\scriptsize] {chuyển hướng};
\draw[arr] (llama) -- (gbnf);

\end{tikzpicture}
\caption{Việc chuyển serving engine bị ép buộc bởi phần cứng, nhưng đó
cũng chính là điều làm cho giải mã ràng buộc theo schema ở
Hình~\ref{fig:constrained-decoding} khả thi trên máy này —
\texttt{src/grader/llm\_client.py} không cần thay đổi gì ngoài
\texttt{base\_url} vì cả hai engine đều phơi bày API tương thích OpenAI.}
\label{fig:pivot}
\end{figure}

% =====================================================================
\section{Kết quả định lượng}

\begin{table}[htbp]
\centering
\begin{tabular}{p{4.6cm}p{5.4cm}p{3.6cm}}
\toprule
\textbf{Thực nghiệm} & \textbf{Thiết lập} & \textbf{Kết quả} \\
\midrule
Consistency, không batching & 4 mẫu $\times$ 100 lần chạy, \texttt{--parallel 1} & exact-match = 1.0, cả 4 mẫu \\
Consistency, có batching & Cùng 4 mẫu $\times$ 100 lần chạy, \texttt{--parallel 4}, 8 đồng thời & exact-match = 1.0, cả 4 mẫu \\
Consistency, mẫu khó & 4 bài nộp có lỗi/trường hợp biên $\times$ 100 lần, cả hai điều kiện & exact-match = 1.0, cả 4 mẫu \\
Ablation (căn cứ kiến trúc) & 65 mẫu, 3 nhánh, Spearman $\rho$ so với nhãn tham chiếu & 0.60 $\rightarrow$ 0.85 $\rightarrow$ \textbf{0.88} \\
Parse fail, các nhánh free-form & 65 mẫu, cả hai điều kiện free-form & 0/65 \\
\bottomrule
\end{tabular}
\caption{Các con số chính, đều lấy từ \texttt{docs/thesis/results.md}
(không tính lại ở đây).}
\label{tab:results}
\end{table}

% =====================================================================
\section{Stack công nghệ}

\begin{figure}[htbp]
\centering
\begin{tikzpicture}[>=Latex]
\def\w{11.5}
\node[iobox, text width=\w cm, fill=black!4] (l5) at (0,4.8) {\textbf{Demo UI} --- \texttt{streamlit} (\texttt{app.py}): lớp vỏ mỏng quanh pipeline thật, không phải bản triển khai lại};
\node[iobox, text width=\w cm, fill=orange!7] (l4) at (0,3.4) {\textbf{Experiments} --- \texttt{scipy} (tương quan Spearman); harness tự viết cho exact-match và ablation 3 nhánh};
\node[iobox, text width=\w cm, fill=blue!6] (l3) at (0,2.0) {\textbf{Grading core} --- Python thuần, \texttt{pydantic} (định nghĩa/kiểm tra schema), SDK \texttt{openai} làm client mỏng};
\node[iobox, text width=\w cm, fill=green!7] (l2) at (0,0.6) {\textbf{Serving} --- \texttt{llama.cpp} \texttt{llama-server}, backend Metal, JSON ràng buộc GBNF, API tương thích OpenAI};
\node[iobox, text width=\w cm, fill=black!7] (l1) at (0,-0.8) {\textbf{Model} --- Qwen2.5-Coder-7B-Instruct (GGUF, Q4\_K\_M) --- model lớn nhất còn nằm trong phạm vi sub-7B của dự án};

\draw[arr] (l5.south) -- (l4.north);
\draw[arr] (l4.south) -- (l3.north);
\draw[arr] (l3.south) -- (l2.north);
\draw[arr] (l2.south) -- (l1.north);
\end{tikzpicture}
\caption{Stack phân lớp, lớp trên sử dụng lớp dưới.}
\label{fig:stack}
\end{figure}

% =====================================================================
\section{Thực hành kỹ thuật đáng chú ý}

\begin{itemize}[leftmargin=1.4em]
  \item \textbf{Chung một đường code giữa bộ sinh dữ liệu và pipeline
    thật:} \texttt{sandbox.run\_submission()} là hàm giống hệt mà
    \texttt{data/build\_dataset.py} dùng để sinh nhãn tham chiếu, nên
    pipeline chấm điểm và bộ dữ liệu dùng để kiểm chứng nó không thể âm
    thầm lệch nhau theo thời gian.
  \item \textbf{Dữ liệu sinh ra, không commit:}
    \texttt{data/dataset.json} được sinh từ \texttt{data/problems.py} qua
    \texttt{data/build\_dataset.py} và bị gitignore — nguồn sự thật là bộ
    sinh, không phải một file JSON cũ có thể lỗi thời.
  \item \textbf{Lưu lại các lần so sánh, không ghi đè kết quả:} mỗi lần
    chạy lại ở một kích thước model hoặc dataset mới đều giữ lại kết quả
    trước đó (\texttt{results/archive\_1.5b/},
    \texttt{results/ablation\_33samples\_7b.json}) để các tuyên bố
    trước/sau có thể kiểm chứng được, không chỉ là khẳng định suông.
  \item \textbf{Hạn chế đã biết được nêu ngay tại nơi dữ liệu sống, không
    giấu đi:} \texttt{data/script.md} ghi rõ, ngay trong thư mục của bộ dữ
    liệu, rằng nhãn \texttt{efficiency}/\texttt{code\_style} là tự gán
    nhãn một người, không phải ground truth độc lập — mọi bảng kết quả
    dùng các nhãn này đều nhắc lại lưu ý đó.
\end{itemize}

% =====================================================================
\section{Những gì cố tình nằm ngoài phạm vi (và vì sao đó là một quyết định)}

\begin{itemize}[leftmargin=1.4em]
  \item \textbf{Fine-tuning (LoRA):} được xây như một phương án dự phòng,
    có cổng go/no-go. Vì prompting + giải mã có cấu trúc đã vượt ngưỡng
    consistency và agreement của dự án (Bảng~\ref{tab:results}), giai đoạn
    fine-tuning bị bỏ qua — ghi nhận như một phát hiện
    (\texttt{docs/roadmap.md}, Phase 4), không phải một tính năng dang dở.
  \item \textbf{Kiểm chứng bằng giám khảo con người thật:} các con số
    agreement hiện tại là so với nhãn tham chiếu dựa trên luật, không bao
    giờ được gọi là ``đồng thuận với giám khảo con người'' ở bất kỳ đâu
    trong tài liệu repo (xem lưu ý trong \texttt{data/script.md}). Việc
    tuyển một giám khảo độc lập là một bước tiếp theo đã được lên kế
    hoạch nhưng chưa bắt đầu (\texttt{docs/roadmap.md} T3.2), tách bạch
    rõ ràng khỏi các con số không phụ thuộc vào nó.
\end{itemize}

\end{document}
